% BHU PYQ -- Professional Exam Template
\documentclass[11pt,a4paper]{article}

\usepackage[a4paper,top=2.5cm,bottom=2.5cm,left=2.8cm,right=2.8cm,headheight=14pt]{geometry}
\usepackage{amsmath,amssymb}
\usepackage[T1]{fontenc}
\usepackage[utf8]{inputenc}
\usepackage{lmodern}
\usepackage{microtype}
\usepackage{fancyhdr}
\usepackage{enumitem}
\usepackage{listings}
\usepackage{hyperref}
\usepackage{lastpage}
\usepackage{parskip}

\hypersetup{colorlinks=true,urlcolor=black,linkcolor=black,
  pdftitle={BVCA-306: GEN -- BHU PYQ}}

\pagestyle{fancy}
\fancyhf{}
\renewcommand{\headrulewidth}{0.4pt}
\renewcommand{\footrulewidth}{0.4pt}
\fancyhead[L]{\small\textit{Banaras Hindu University}}
\fancyhead[R]{\small\textit{BVCA-306: GEN}}
\fancyfoot[C]{\small Page \thepage\ of \pageref{LastPage}}

\lstset{
  basicstyle=\small\ttfamily,
  frame=single,
  breaklines=true,
  columns=flexible,
  keepspaces=true
}

\newlist{parts}{enumerate}{2}
\setlist[parts,1]{label=(\alph*),leftmargin=*,itemsep=4pt,topsep=3pt}
\setlist[parts,2]{label=(\roman*),leftmargin=*,itemsep=2pt,topsep=2pt}
\newcommand{\pts}[1]{\hfill\textit{[#1]}}

\begin{document}
\begin{center}
  {\large\bfseries BANARAS HINDU UNIVERSITY}\\[2pt]
  {\normalsize B.Voc. Semester III Examination, 2023-24}\\[6pt]
  \rule{0.8\linewidth}{0.5pt}\\[6pt]
  {\large\bfseries Paper: BVCA-306 --- GEN}\\[4pt]
  {\normalsize Time: Three hours \hfill Maximum Marks: 70}
\end{center}

\rule{\linewidth}{0.4pt}

\smallskip
\noindent\textit{Instructions: 1, Attempt any two out of four questions in Section A.}

\bigskip

GEN - 302 : Introduction to IoT
Instructions:
1, Attempt any two out of four questions in Section A. (2*12=24 Marks)

\medskip\noindent\textbf{Question 2.} Attempt any six out of nine questions in Section B (6*6=36 Marks)

\medskip\noindent\textbf{Question 3.} Attempt all 10 questions in Section C (1*10=10 Marks )

\bigskip\noindent{\large\bfseries SECTION --- A}
\rule{\linewidth}{0.4pt}\smallskip

\medskip\noindent\textbf{Question 1.} Explain 3 and 5-layer IoT architecture in detail.

\medskip\noindent\textbf{Question 2.} Differentiate between microprocessors and microcontrollers. What are the
applications areas of Microprocessors and Microcontrollers?

\medskip\noindent\textbf{Question 3.} What is the purpose of sensors in IoT? What are the required criteria to choose a
suitable sensor to measure the desired physical parameter?

\medskip\noindent\textbf{Question 4.} What do you understand by the term loT Ecosystem? What are the components
of an IoT system? Describe all.

\bigskip\noindent{\large\bfseries SECTION --- B :}
\rule{\linewidth}{0.4pt}\smallskip

1, What do you mean by IoT? What are the advantages and disadvantages of IoT
applications?

\medskip\noindent\textbf{Question 2.} What is cloud? What are the basic characteristics of cloud computing?

\medskip\noindent\textbf{Question 3.} Differentiate RISC and CISC microprocessors.

\medskip\noindent\textbf{Question 4.} What is the difference between transducers, sensors and actuators?

\medskip\noindent\textbf{Question 5.} What do you mean by Embedded System? Explain its characteristics.

\medskip\noindent\textbf{Question 6.} Explain the components of the Microcontroller.

\medskip\noindent\textbf{Question 7.} Define the Term: a) Data Processing b) Automation c) Communication Protocols

\medskip\noindent\textbf{Question 8.} Explain any 6 applications of IoT.

\medskip\noindent\textbf{Question 9.} Discuss the principles of operation for each type of sensor and their applications
in IoT devices.

\bigskip\noindent{\large\bfseries SECTION --- C}
\rule{\linewidth}{0.4pt}\smallskip

\medskip\noindent\textbf{Question 1.} Which protocol provide Full-Duplex communication?
\begin{parts}
  \item WebSocket
  \item HTTP
  \item MQTT
  \item CoAP
\end{parts}

\medskip\noindent\textbf{Question 2.} The RFID technology consists of:
\begin{parts}
  \item Antenna
  \item Transponder
  \item Transceiver
  -d. All of these
  Fr LsO:
  Q)
\end{parts}

\medskip\noindent\textbf{Question 3.} Connected car is :
\begin{parts}
  \item BAN
  B.LAN
  C. WAN
  D. None Of These
\end{parts}

\medskip\noindent\textbf{Question 4.} The full form of LoRa is: :
\begin{parts}
  \item label range radio |
  \item long radio range ,
  \item none of these
  \_5. The different wireless networks used in cellular network: }
  \item GSM : |
  B, CDMA
  C. LTE i
  D. All Of The Above
\end{parts}

\medskip\noindent\textbf{Question 6.} Which of the following is Lightweight protocol used in IoT Applications?
\begin{parts}
  \item MQTT
  : b. WebSocket
  \item HTTP ,
  d All of the above 3
\end{parts}

\medskip\noindent\textbf{Question 7.} Which of the following wireless technologies provide mesh network architecture in
easiest way? : |
\begin{parts}
  \item Bluetooth |
  \item ZigBee
  \item Wi-Fi
  \item None of these
\end{parts}

\medskip\noindent\textbf{Question 8.} What is the full form of HTTP?
\begin{parts}
  \item Hyphenation Text Test Program
  \item Hypertext Transfer Protocol
  \item Hypertext Transfer Package
  \item None of the above
\end{parts}

\medskip\noindent\textbf{Question 9.} An entity in IoT:
\begin{parts}
  \item Businesses
  \item Governments
  \item Consumers
  \item All of the above
\end{parts}

\medskip\noindent\textbf{Question 10.} ZigBee is based on which of the following standard:
\begin{parts}
  \item IEEE 802.15.4
  B. IEEE 802.5
  C. IEEE 802.4.13
  D. Both a and b
  i 2 oo ee eo ok ok 2 ok kK oe
\end{parts}

\vfill
\rule{\linewidth}{0.4pt}
\begin{center}\small
\href{https://bhu-exam.vercel.app/}{Click here to visit our website}\\[4pt]\href{https://me-aryan.vercel.app/}{Click here to visit our contributor}\\[4pt]\href{https://quashan-me.vercel.app/}{Click here to visit our contributor}
\end{center}
\end{document}